% =====================================================================================================
\section{Cell layout and the four-dimensional keyhole comparator}\label{sec:cmp-4d}

\subsection{The flat model, fundamental domain, and cell embeddings $C_i$}\label{sec:cmp-cells}

Let $k\ge10$ and $\eta\in(0,1-\sqrt{99/100}]$, and set $s=\sqrt k>3$, $t_0=1-\frac12\eta$, $a=t_0/s$, $h=t_0/3$,
$c_s=3s-9>0$, $g=3(1-t_0)=\frac32\eta$, and $T=g/c_s$. On the flat chart $(G,v,w)\in\R\times\R\times\C$ with symplectic
form $\Omega_{\mathrm{fl}}=dG\wedge dv+dw_x\wedge dw_y$, define $P(p)=3-9p$, $J(p)=3sp-3$, and $K(p)=c_sp$ for
$p=\pi|w|^2<\frac13$. Notice the exact period identity
\begin{equation}\label{eq:JK-period}
 K(p)-J(p) = (3s-9)p - (3sp-3) = 3-9p = P(p) > 0,
\end{equation}
and the deck group generated by the two $C^\omega$ symplectomorphisms
\[
 T_u(G,v,w)=\bigl(G+P(p),\,v,\,we^{2\pi i(9v)}\bigr),\qquad T_v(G,v,w)=(G,v+1,w),
\]
which commute modulo integer fibre turns ($T_uT_v=T_vT_u$ because $e^{2\pi i\cdot9}=1$) and have fundamental domain
$\mathcal F_{\mathrm{fl}}=\{p<\frac13,\ J(p)\le G<K(p),\ 0\le v<1\}$. Define the level curves
\[
 O(G)=\frac{G+3}{3s},\quad E_{\mathrm{lin}}(G)=\frac{G}{c_s},\quad
 P_{\mathrm{top}}(G)=\frac{ka-3-G}{k-3s},\quad \operatorname{slack}(G)=T\Bigl(1+\frac{\min(G,0)}{g}\Bigr).
\]
Direct algebra gives the linear comparison identity of Paper~I in level-$G$ form:
\begin{equation}\label{eq:level-G-id}
 O(G)-E_{\mathrm{lin}}(G)-P_{\mathrm{top}}(G) \equiv T \qquad(\forall\,G\in\R),
\end{equation}
and $p<O(G)\iff G>J(p)$, while $p>E_{\mathrm{lin}}(G)\iff G<K(p)$.

\begin{lemma}[Cell-to-flat embeddings $C_i$]\label{lem:cell-flat}
For $0\le i\le k-1$, define $C_i\colon\mathcal C_a\to\R\times\R\times\C$ on $(X,T_{\mathrm{cell}},w)$ ($p=\pi|w|^2<a$,
$0<X<a-p$, $0<T_{\mathrm{cell}}<1$) by
\[
 C_i(X,T_{\mathrm{cell}},w)=
 \begin{cases}
  \bigl(J(p)+9X,\ \frac{i-4+T_{\mathrm{cell}}}{9},\ we^{-2\pi i(3sT_{\mathrm{cell}}/9)}\bigr), & 0\le i<9,\\[4pt]
  \bigl(J(p)+i(a-p)+X,\ T_{\mathrm{cell}}+\frac12,\ we^{-2\pi i((3s-i)T_{\mathrm{cell}})}\bigr), & 9\le i\le k-1.
 \end{cases}
\]
Then each $C_i$ is a real-analytic symplectic embedding ($C_i^*\Omega_{\mathrm{fl}}=dX\wedge dT_{\mathrm{cell}}+\omega_V$).
The nine group images $\mathcal G_i=C_i(\mathcal C_a)$ ($0\le i<9$) are pairwise disjoint in $\mathcal F_{\mathrm{fl}}$
(lying in disjoint $v$-windows $(\frac{i-4}{9},\frac{i-3}{9})\subset(-\frac49,\frac59)$) and satisfy
\begin{equation}\label{eq:group-G-bound}
 J(p) < G < J(p)+9(a-p) = (3s-9)p + 9a - 3 < 3sa - 3 = -g\qquad(\forall\,(G,v,w)\in\mathcal G_i).
\end{equation}
The stack images $\mathcal S_i=C_i(\mathcal C_a)$ ($9\le i<k$) are pairwise disjoint ( separated by the hyperplanes
$A_i(p)=J(p)+i(a-p)$), satisfy $G>J(p)+9(a-p)$, $G<J(0)+ka=st_0-3=:G_{\max}$, and at each level $G$ have fibre slice
$p\le P_{\mathrm{top}}(G)$ (with $p<a$).
\end{lemma}

\begin{proof}
For $i<9$, $d(J(p)+9X)\wedge d(\frac{i-4+T_{\mathrm{cell}}}{9})=(3s\,dp+9\,dX)\wedge\frac{dT_{\mathrm{cell}}}{9}
=dX\wedge dT_{\mathrm{cell}}+\frac{s}{3}dp\wedge dT_{\mathrm{cell}}$, which cancels $dp\wedge d(\theta-\frac{s}{3}T_{\mathrm{cell}})-dp\wedge d\theta$;
similarly for $i\ge9$, $d(J(p)+i(a-p)+X)\wedge dT_{\mathrm{cell}}=dX\wedge dT_{\mathrm{cell}}+(3s-i)dp\wedge dT_{\mathrm{cell}}$.
Since $3s-9>0$ and $p<a$, $(3s-9)p+9a-3<(3s-9)a+9a-3=3sa-3=3t_0-3=-g$. Finally, $G<J(p)+k(a-p)=-3+ka-(k-3s)p$ is
equivalent to $p<P_{\mathrm{top}}(G)$.
\end{proof}

\subsection{The four-dimensional keyhole comparator $F=F_2\circ F_1$}\label{sec:cmp-F}

Fix $0<\lambda_2<\lambda_a<1$ and $\mu\in(0,1)$ ($\lambda_a=\frac9{10}$, $\lambda_2=\frac3{10}$, $\mu=\frac12$), and
put $w_G=\lambda_a g$ and $w_2=\lambda_2 g$. Let $K_C$ be a smooth keyhole on $\bar D(a)$ satisfying \eqref{eq:K123}
with hole $x_0=\pi r_h^2$ and slit half-width $\theta_0=W_1$. Define
\[
 \tau(G)=S_4\Bigl(\frac{G+w_G}{w_G-w_2}\Bigr),\quad
 c(G)=\frac{w_2}{c_s}R_4\Bigl(\frac{G+w_2}{w_2}\Bigr),\quad
 \vartheta(\theta)=\theta-\tfrac12\sgn(\theta)\ \ \bigl(\theta\in(-\tfrac12,\tfrac12]\setminus\{0\}\bigr).
\]
Notice that $\vartheta(\arg(w)/(2\pi))=\arg(-w)/(2\pi)\in(-\frac12,\frac12)$ is real-analytic on the slit plane
$\C\setminus[0,\infty)$ (in particular across the negative real axis $\theta=\pm\frac12$!). Define $F_1(G,v,w)=\bigl(G,v-\tau'(G)g_{\tau(G)}(w),KS_{\tau(G)}(w)\bigr)$ and, on $\{c(G)=0\}\cup\{(G,v,w_1):w_1\in\C\setminus[0,\infty),\ \pi|w_1|^2+c(G)>0\}$, define $F_2(G,v_1,w_1)=(G,v_1,w_1)$ when $c(G)=0$ ($G\le-w_2$) and
\[
 F_2(G,v_1,w_1)=\Bigl(G,\ v_1-c'(G)\vartheta(\theta_1),\ \sqrt{\frac{x_1+c(G)}{\pi}}\,e^{2\pi i\theta_1}\Bigr)
 \qquad(x_1=\pi|w_1|^2,\ \theta_1=\arg(w_1)/(2\pi))
\]
when $c(G)>0$ ($G>-w_2$). Set $F=F_2\circ F_1$, $F^{(i)}=\mathrm{id}$ for $0\le i<9$, and $F^{(i)}=F$ for $9\le i<k$.

\begin{theorem}[Domain, symplecticity, and stack-corridor separation of $F$]\label{thm:comparator}
Suppose $K_C$ satisfies \eqref{eq:K123} on $\bar D(a)$ together with the schedule inequalities:
\begin{itemize}
\item \emph{(H6)} $\max\bigl(S_4(u)\delta_{\mathrm{rel}},\delta_{\mathrm{small}}\bigr)\le(1-\mu)T\bigl(1-\lambda_a+u(\lambda_a-\lambda_2)\bigr)$ for $u\in[0,1]$;
\item \emph{(H7)} $\max(\delta_{\mathrm{rel}},\delta_{\mathrm{small}})+\lambda_2 T R_4(u)\le(1-\mu)T\bigl(1-\lambda_2(1-u)\bigr)$ for $u\in[0,1]$;
\item \emph{(H8)} $\pi d^2\le T(1-\mu-\frac12\lambda_2)$ and $\pi d^2+\frac12\lambda_2 T+\mu T\le O(-w_G)$.
\end{itemize}
Then on the open stack union $\mathcal S=\bigcup_{i=9}^{k-1}\mathcal S_i$:
\begin{enumerate}[label=(\roman*)]
\item For every $(G,v,w)\in\mathcal S$ with $G>-w_2$, $\tau(G)=1$ and $w\in\Bo(a)$, so by Theorem~\ref{thm:kh-bounds}(ii)
$w_1=KS_1(w)$ satisfies $x_1>x_0>0$ and $|\theta_1|>\theta_0>0$, avoiding $[0,\infty)$ by an open sector. As $G\downarrow-w_2$,
$c(G)$ and $c'(G)$ vanish to order $9$ while $|\vartheta(\theta_1)|<\frac12$ and $x_1>x_0$. Hence $F$ is a $C^7$
symplectic embedding on an open neighbourhood of $\mathcal S$, with explicit $C^7$ inverse $F_1^{-1}\circ F_2^{-1}$.
\item Every image point $(G,v_2,w_2)\in F(\mathcal S)$ ($p_2=\pi|w_2|^2$) lies strictly inside a single fundamental
domain $[J(p_2),K(p_2))$ of $T_u$ and satisfies:
\begin{align}
 p_2 &< h_K \qquad\text{(whenever $h_K\ge a$ and $h_K>\max_{G\in[-w_G,G_{\max}]}(O(G)-\mu\operatorname{slack}(G))$)},
 \label{eq:p2-hK}\\
 J(p_2) &< G < K(p_2) - \min\bigl(w_2,\ c_s x_0+\tfrac12 w_2\bigr) \le K(p_2)-\tfrac12 w_2,
 \label{eq:stack-upper}\\
 G &> J(p_2)+9(a-p_2)\qquad\text{whenever }p_2<a.\label{eq:stack-above-group}
\end{align}
Consequently, the $k$ sets $\mathcal G_0,\dots,\mathcal G_8,F(\mathcal S_9),\dots,F(\mathcal S_{k-1})$ are pairwise
disjoint modulo the deck group $\langle T_u,T_v\rangle$, and all of them miss the corridor
$\mathcal K_{\mathrm{free}}=\{p<\frac13:\ K(p)-\frac12w_2\le G<K(p)\}\pmod{P(p)}$.
\end{enumerate}
\end{theorem}

\begin{proof}
(i) $F_1$ is a global $C^7$ symplectomorphism by Proposition~\ref{prop:ks-symp} ($\tau'\in C^7$). On $\{G>-w_2\}$,
$d(x_1+c(G))\wedge d\theta_1 + dG\wedge d(v_1-c'(G)\vartheta(\theta_1))
=dx_1\wedge d\theta_1+c'(G)dG\wedge d\theta_1+dG\wedge dv_1-c'(G)dG\wedge d\vartheta(\theta_1)=dG\wedge dv_1+dx_1\wedge d\theta_1$
since $\vartheta'(\theta_1)=1$ off $\theta_1=0$. Because $\mathcal S_9,\dots,\mathcal S_{k-1}$ are pairwise disjoint and
$F$ is injective on their union, $F(\mathcal S_9),\dots,F(\mathcal S_{k-1})$ are pairwise disjoint.

(ii) If $G\le-w_G$, $F=\mathrm{id}$, so $p_2=p<a\le h_K$, $J(p_2)+9(a-p_2)<G\le-w_G<-w_2<K(p_2)-w_2$.
Now suppose $G\in(-w_G,G_{\max})$. Since $w\in\Bo(\min(a,P_{\mathrm{top}}(G)))$, Theorem~\ref{thm:kh-bounds}(i) gives
$x_1\le\max(P_{\mathrm{top}}(G)+\max(\tau(G)\delta_{\mathrm{rel}},\delta_{\mathrm{small}}),\pi d^2)$, and adding $c(G)$
gives $p_2=x_1+c(G)$.
On $(-w_G,-w_2]$ ($c(G)=0$), (H6) gives $P_{\mathrm{top}}(G)+\max(\tau\delta_{\mathrm{rel}},\delta_{\mathrm{small}})
\le P_{\mathrm{top}}(G)+(1-\mu)\operatorname{slack}(G)$. Notice that for $G\le0$, $G/c_s=T(G/g)$, so \eqref{eq:level-G-id}
gives the exact identity
\[
 O(G)-P_{\mathrm{top}}(G) = T + \frac{G}{c_s} = T\Bigl(1+\frac{G}{g}\Bigr) = \operatorname{slack}(G)\qquad(G\le0),
\]
whereas for $G\ge0$, $O(G)-P_{\mathrm{top}}(G)=T+E_{\mathrm{lin}}(G)$ with $\operatorname{slack}(G)=T$. Therefore:
\begin{itemize}
\item On $(-w_G,-w_2]$, (H6) gives $P_{\mathrm{top}}(G)+\max(\tau\delta_{\mathrm{rel}},\delta_{\mathrm{small}})\le O(G)-\mu\operatorname{slack}(G)$;
\item On $[-w_2,0]$, (H7) gives $P_{\mathrm{top}}(G)+\max(\delta_{\mathrm{rel}},\delta_{\mathrm{small}})+c(G)\le O(G)-\mu\operatorname{slack}(G)$;
\item On $[0,G_{\max}]$, $c(G)=\frac{G}{c_s}+\frac12\lambda_2 T$, and (H7) at $u=1$ gives
$\max(\delta_{\mathrm{rel}},\delta_{\mathrm{small}})+\frac12\lambda_2 T\le(1-\mu)T$, so
$P_{\mathrm{top}}(G)+\max(\delta_{\mathrm{rel}},\delta_{\mathrm{small}})+c(G)\le P_{\mathrm{top}}(G)+\frac{G}{c_s}+(1-\mu)T=O(G)-\mu T=O(G)-\mu\operatorname{slack}(G)$.
\end{itemize}
In the remaining $\pi d^2+c(G)$ branch of \eqref{eq:kh-nest}: on $(-w_G,0]$, $c(G)\le c(0)=\frac12\lambda_2 T$ and
$\operatorname{slack}(G)\le T$, so (H8) gives $\pi d^2+c(G)+\mu\operatorname{slack}(G)\le\pi d^2+\frac12\lambda_2 T+\mu T\le O(-w_G)<O(G)$;
on $[0,G_{\max}]$, $O(G)-\frac{G}{c_s}=P_{\mathrm{top}}(G)+T\ge T$, so (H8) gives
$\pi d^2+c(G)+\mu T=\pi d^2+\frac{G}{c_s}+(\frac12\lambda_2+\mu)T\le\frac{G}{c_s}+T\le O(G)$.
Thus $p_2\le O(G)-\mu\operatorname{slack}(G)<O(G)$ holds for \emph{all} $G\in(-w_G,G_{\max})$, which gives both
\eqref{eq:p2-hK} and $G>3sp_2-3=J(p_2)$!

For the upper bound \eqref{eq:stack-upper}: if $G\in(-w_G,-w_2]$, then $K(p_2)\ge0$, so $K(p_2)-G\ge-G\ge w_2$. If
$G\in(-w_2,G_{\max})$, Theorem~\ref{thm:kh-bounds}(ii) gives $x_1>x_0$, so $p_2>x_0+c(G)$; writing $u=(G+w_2)/w_2>0$
($G=(u-1)w_2$) and using $u-R_4(u)\le\lim_{t\to\infty}(t-R_4(t))=\frac12$, we have
\[
 K(p_2)-G > c_sx_0 + c_sc(G) - G = c_sx_0 + w_2\bigl(1-(u-R_4(u))\bigr) \ge c_sx_0 + \tfrac12 w_2.
\]
Finally, for $G>-w_G=-0.9g$ and any $p_2<a$, \eqref{eq:group-G-bound} gives $J(p_2)+9(a-p_2)<-g<-0.9g<G$, proving
\eqref{eq:stack-above-group} and disjointness modulo $\langle T_u,T_v\rangle$.
\end{proof}

% =====================================================================================================
\section{Marked-fibre twists and corridor routing}\label{sec:twists}

In the flat chart, put $u=G/P(p)$ ($P(p)=3-9p>0$). The three straight marked fibres of the cubic transfer (Section~\ref{sec:transfer})
are $\mathcal M_1=\{(u,v)\equiv(\frac56,0)\}$, $\mathcal M_2=\{(u,v)\equiv(\frac13,0)\}$, and
$\mathcal M_3=\{(u,v)\equiv(0,\frac12)\}$ modulo $\Z^2$. Since the ball images fill $(u,v)$ across $(u_j,v_j)$, we
interpose two closed-form $C^4$ deck-equivariant symplectomorphisms $\Psi=\Phi_g\circ\Psi_S$ that pull all three
marked fibres back into the empty corridor $\mathcal K_{\mathrm{free}}$.

\begin{construction}[Toric twists $\Psi_S$ and $\Phi_g$]\label{cons:twists}
Fix constants $a_g=\frac1{10}$, $(\omega_{m,0},\omega_{m,1})=(0.015,0.045)$, $(\omega_{u,0},\omega_{u,1})=(0.02,0.12)$,
$v_{\mathrm{comp}}=\frac14$, $(\omega_{c,0},\omega_{c,1})=(0.05,0.14)$, $\delta=\frac14w_2$, $u_*=-c_s/9$,
$c_0=\delta-c_s/3$, $n_j=-\operatorname{round}(u_j-u_*)\in\Z$, $\phi_j=u_j+n_j-u_*$, and shifted $v$-targets
$(v_1',v_2',v_3')=(-a_g,0,\frac12)$. Using $m=2$ plateau bumps $\beta_v=\beta_{\omega_{m,0},\omega_{m,1},2}$,
$\beta_c=\beta_{\omega_{c,0},\omega_{c,1},2}$ (mass $M_c=\omega_{c,0}+\omega_{c,1}$), and
$\beta_u=\beta_{\omega_{u,0},\omega_{u,1},2}$:
\begin{enumerate}[label=(\arabic*)]
\item For $v\in\R$, reduce $\tilde v\in[-\frac13,\frac23)$ modulo $1$ and set
$b_j(v)=\beta_v(\tilde v-v_j')-\frac{M_v}{M_c}\beta_c(\tilde v-v_{\mathrm{comp}})$,
$B_j(v)=\mathcal B_v(\tilde v-v_j')-\frac{M_v}{M_c}\mathcal B_c(\tilde v-v_{\mathrm{comp}})$ (so $B_j(v)$ is
$1$-periodic in $C^5(\R)$ because $\int_{-1/3}^{2/3}b_j=0$). Put
$S_v(p,v)=c_0+P(p)\sum_{j=1}^3\phi_j b_j(v)$, $S_p(v)=-9\sum_{j=1}^3\phi_j B_j(v)$, and
\[
 \Psi_S(G,v,w)=\bigl(G+S_v(p,v),\ v,\ we^{-2\pi i S_p(v)}\bigr).
\]
\item For $u\in\R$, reduce $\tilde u\in[\frac15,\frac65)$ modulo $1$ and set
$g_{\mathrm{tw}}(u)=-a_g\bigl(\beta_u(\tilde u-\frac56)-\beta_u(\tilde u-\frac7{12})\bigr)$ and
$G_g(u)=-a_g\bigl(\mathcal B_u(\tilde u-\frac56)-\mathcal B_u(\tilde u-\frac7{12})\bigr)$ ($1$-periodic in $C^4(\R)$
and $C^5(\R)$). Define
\[
 \Phi_g(G,v,w)=\bigl(G,\ v-g_{\mathrm{tw}}(u),\ we^{2\pi i\cdot9(G_g(u)-ug_{\mathrm{tw}}(u))}\bigr)\qquad(u=G/P(p)).
\]
\end{enumerate}
\end{construction}

\begin{proposition}[Symplecticity, deck equivariance, and marked-fibre routing]\label{prop:twists}
On $\{p<\frac13\}$:
\begin{enumerate}[label=(\roman*)]
\item $\Psi_S$ and $\Phi_g$ are $C^4$ symplectomorphisms preserving $p=\pi|w|^2$, with closed-form inverses obtained by
reversing the sign of $(S_v,S_p)$ and $g_{\mathrm{tw}}$, and both commute with $T_u$ and $T_v$.
\item Under the window-separation checks \emph{(H10)} (the four $v$-intervals $[v_j'\pm0.045]$, $[0.25\pm0.14]$ are
pairwise disjoint in $[-\frac13,\frac23)$, and the two $u$-lobes $[\frac56\pm0.12]$, $[\frac7{12}\pm0.12]$ are disjoint
in $[\frac15,\frac65)$ and miss $\frac13$ and $1$), the preimage of the deck orbit of $\mathcal M_j$ under
$\Psi=\Phi_g\circ\Psi_S$ is the exact spiral
\begin{equation}\label{eq:spiral}
 \Psi^{-1}\bigl(\operatorname{deck}(\mathcal M_j)\bigr)
 = \bigl\{(G,v,w): G\equiv K(p)-\delta\pmod{P(p)},\ v\equiv v_j'\pmod1\bigr\}.
\end{equation}
Since $0<\delta=\frac14w_2<\frac12w_2$ and $3-9a>\delta$ \emph{(H9b)}, all three spirals lie inside the free corridor
$\mathcal K_{\mathrm{free}}$ and are disjoint from $\mathcal G_0,\dots,\mathcal G_8,F(\mathcal S_9),\dots,F(\mathcal S_{k-1})$.
\end{enumerate}
\end{proposition}

\begin{proof}
(i) Writing $S(p,v)=c_0v+P(p)\sum_j\phi_jB_j(v)$ locally, $\partial_vS=S_v$ and $\partial_pS=S_p$, so
$d(G+S_v)\wedge dv+dp\wedge d(\theta-S_p)=dG\wedge dv+dp\wedge d\theta$. For $\Phi_g$, writing $H_g(G,p)=P(p)G_g(G/P(p))$
gives $\partial_GH_g=g_{\mathrm{tw}}(u)$ and $\partial_pH_g=-9(G_g(u)-ug_{\mathrm{tw}}(u))$, so $\Phi_g$ is the exact
time-one flow of $H_g$ (constant along orbits since $G$ and $p$ are preserved) and is $C^4$ because $g_{\mathrm{tw}}\in C^4$.
Deck equivariance holds because $b_j,B_j$ are $1$-periodic in $v$ and $g_{\mathrm{tw}},G_g$ are $1$-periodic in $u$
(for $T_u$ on $\Phi_g$: shifting $u\mapsto u+1$ shifts $u g_{\mathrm{tw}}(u)$ by $g_{\mathrm{tw}}(u)$, shifting the
fibre phase by $-9g_{\mathrm{tw}}(u)$, which matches $T_u$'s phase $9(v-g_{\mathrm{tw}}(u))$).

(ii) At $u_1\equiv\frac56$, $g_{\mathrm{tw}}(\frac56)=-a_g(1-0)=-a_g$, whereas at $u_2\equiv\frac13$ and $u_3\equiv0$,
$g_{\mathrm{tw}}=0$ by (H10). Thus $\Phi_g^{-1}(\operatorname{deck}(\mathcal M_j))=\{u\equiv u_j\pmod1,\ v\equiv v_j'\pmod1\}$.
At $v\equiv v_j'\pmod1$, (H10) gives $b_j(v_j')=1$, $b_l(v_j')=0$ ($l\ne j$), and $\beta_c(v_j'-v_{\mathrm{comp}})=0$,
so $S_v(p,v_j')=c_0+P(p)\phi_j$. Since $\phi_j=u_j+n_j+c_s/9$ and $c_0=\delta-c_s/3$,
\[
 S_v(p,v_j') = (3-9p)\Bigl(u_j+n_j+\frac{c_s}{9}\Bigr)+\delta-\frac{c_s}{3}
 = P(p)(u_j+n_j) - c_sp + \delta = P(p)(u_j+n_j) - \bigl(K(p)-\delta\bigr).
\]
Hence $G+S_v(p,v_j')\equiv P(p)u_j\pmod{P(p)}$ if and only if $G\equiv K(p)-\delta\pmod{P(p)}$, proving \eqref{eq:spiral}.
Since $0<\delta=\frac14w_2<\min(\frac12w_2,3-9a)$, in the fundamental domain $[J(p),K(p))$ we have
$J(p)+9(a-p)=K(p)-(3-9a)<K(p)-\delta$ (for $p<a$) and $\sup F(\mathcal S)_G<K(p)-\frac12w_2<K(p)-\delta$ by
\eqref{eq:stack-upper}.
\end{proof}
